% =================================================================
%  style/theorem-boxes.tex
%  Цветные theorem-like боксы конспекта (tcolorbox)
%
%  Синтаксис у всех боксов одинаков:
%      \begin{Theorem}{<заголовок или пусто>}{<метка или пусто>}
%          ...
%      \end{Theorem}
%  Необязательный первый аргумент — дополнительные опции tcolorbox,
%  через него подключается пометка обязательного пункта:
%      \begin{Definition}[required]{Линейный оператор}{}
%  Ссылка на бокс с меткой foo: \ref{thm:foo} (префикс зависит от бокса:
%  thm, lem, prop, stmt, def, ex, rem, cor).
% =================================================================

% xcolor уже загружен в style/preamble.tex с опцией [dvipsnames];
% здесь он нужен только как зависимость.
\usepackage{xcolor}

\usepackage[most]{tcolorbox}
% breakable оставлен загруженным на случай, если какому-то боксу
% всё же понадобится разрыв: ему хватит локального breakable.
\tcbuselibrary{breakable,skins,theorems}
\usepackage{enumitem}

% -----------------------------------------------------------------
% Ссылки
% -----------------------------------------------------------------
\definecolor{linkcolor}{HTML}{47528F}

\usepackage{hyperref}
\hypersetup{
	colorlinks=true,
	linkcolor=linkcolor,
	urlcolor=linkcolor,
	citecolor=linkcolor,
	linktoc=section
}

% -----------------------------------------------------------------
% Палитра боксов: <имя>Frame — рамка и плашка заголовка,
%                 <имя>Back  — исходный «бумажный» фон (сейчас
%                 фон боксов строится как Frame!7!white)
% -----------------------------------------------------------------
\definecolor{thmBack}{HTML}{FCF8F9}   \definecolor{thmFrame}{HTML}{8D4455}  % Теорема
\definecolor{lemBack}{HTML}{FCF8FA}   \definecolor{lemFrame}{HTML}{B56C89}  % Лемма
\definecolor{stmtBack}{HTML}{F7FAFD}  \definecolor{stmtFrame}{HTML}{5279A5} % Утверждение
\definecolor{exBack}{HTML}{F7FBF8}    \definecolor{exFrame}{HTML}{5A8F75}   % Пример
\definecolor{defBack}{HTML}{FFFDF8}   \definecolor{defFrame}{HTML}{B68A2E}  % Определение
\definecolor{propBack}{HTML}{FAF8FE}  \definecolor{propFrame}{HTML}{7769A8} % Предложение
\definecolor{remBack}{HTML}{FAFAFA}   \definecolor{remFrame}{HTML}{7E7E7E}  % Замечание
\definecolor{corBack}{HTML}{FDF8F8}   \definecolor{corFrame}{HTML}{C48383}  % Следствие

% -----------------------------------------------------------------
% Базовый стиль бокса
% -----------------------------------------------------------------
\tcbset{
	conspectbox/.style={
		enhanced,
		unbreakable,
		colback=white,
		boxrule=0.8pt,
		arc=2.8mm,
		outer arc=2.8mm,
		left=8pt,
		right=8pt,
		top=6pt,
		bottom=6pt,
		before skip=10pt,
		after skip=10pt,
		fonttitle=\bfseries,
		coltitle=white,
		attach boxed title to top left={
			xshift=6pt,
			yshift*=-\tcboxedtitleheight/2
		},
		boxed title style={
			boxrule=0pt,
			arc=3mm,
			outer arc=3mm,
			left=7pt,
			right=7pt,
			top=2.3pt,
			bottom=2.3pt,
		},
	}
}

% -----------------------------------------------------------------
% Вспомогательный макрос: новый бокс в общем стиле.
%   #1 — опции счётчика, #2 — имя окружения, #3 — подпись,
%   #4 — цвет, #5 — префикс метки
% -----------------------------------------------------------------
% leftrule=1.9pt даёт утолщённую левую границу, которая идёт по самой
% рамке и потому корректно вписывается в скругления. borderline west
% для этого не годится: он рисует прямую вертикаль во всю высоту
% ограничивающего прямоугольника и торчит за скруглённые концы.
\newcommand{\newconspectbox}[5][auto counter, number within=section]{%
	\newtcbtheorem[#1]{#2}{#3}{
		conspectbox,
		colback=#4!7!white,
		colframe=#4,
		colbacktitle=#4,
		leftrule=1.9pt
	}{#5}%
}

% -----------------------------------------------------------------
% Нумерованные боксы (сквозная нумерация внутри section)
% -----------------------------------------------------------------
\newconspectbox{Theorem}{Теорема}{thmFrame}{thm}
\newconspectbox{Lemma}{Лемма}{lemFrame}{lem}
\newconspectbox{Proposition}{Предложение}{propFrame}{prop}
\newconspectbox{Statement}{Утверждение}{stmtFrame}{stmt}

% -----------------------------------------------------------------
% Ненумерованные боксы
% -----------------------------------------------------------------
\newconspectbox[no counter]{Definition}{Определение}{defFrame}{def}
\newconspectbox[no counter]{Example}{Пример}{exFrame}{ex}
\newconspectbox[no counter]{Remark}{Замечание}{remFrame}{rem}
\newconspectbox[no counter]{Corollary}{Следствие}{corFrame}{cor}

% -----------------------------------------------------------------
% Пометка «обязательный пункт»: утолщённая красная левая граница
% и красное (!) на уровне заголовка.
% Подключается через необязательный аргумент бокса:
%     \begin{Definition}[required]{Линейный оператор}{}
%     \begin{Theorem}[required]{Формула Стирлинга}{stirling}
% -----------------------------------------------------------------
\definecolor{reqcolor}{HTML}{C81E2E}

% Красная полоса рисуется поверх левой границы и обрезается по
% скруглённому контуру рамки (arc=2.8mm из conspectbox) — иначе
% прямая вертикаль торчала бы за скруглённые концы.
\newcommand{\conspectreqband}{%
	\begin{scope}
		\clip[rounded corners=2.8mm] (frame.south west) rectangle (frame.north east);
		\fill[reqcolor] (frame.south west) rectangle ([xshift=1.9pt]frame.north west);
	\end{scope}%
}

\tcbset{
	required/.style={
		overlay={
			\conspectreqband
			\node[
			anchor=north east,
			inner sep=0pt,
			font=\bfseries\large,
			text=reqcolor
			] at ([xshift=-8pt, yshift=-12pt]frame.north west) {(!)};
		},
	}
}

% -----------------------------------------------------------------
% Как добавить свой бокс, например «Приём» цвета stmtFrame:
%   \newconspectbox[no counter]{Trick}{Приём}{stmtFrame}{trick}
% -----------------------------------------------------------------
